\documentclass[reqno]{amsart}
\usepackage[T1]{fontenc}
\usepackage{amsmath,amssymb,amsthm}
\usepackage{xr-hyper}
\usepackage[hidelinks]{hyperref}
\externaldocument[d4-]{20261126T060100--a-hecke-operator-is-nearly-onto__workbook}

\newcounter{chap}
% Chapter numbers, generated from 20261125T012000--grh-workbook__guide.org
% by scripts/make-chapters.py.  Do not edit: change the order of the
% index and rerun the script (./preview-site.sh and CI do this).
%
% Usage in a sheet, after \newcounter{chap}:
%   % Chapter numbers, generated from 20261125T012000--grh-workbook__guide.org
% by scripts/make-chapters.py.  Do not edit: change the order of the
% index and rerun the script (./preview-site.sh and CI do this).
%
% Usage in a sheet, after \newcounter{chap}:
%   % Chapter numbers, generated from 20261125T012000--grh-workbook__guide.org
% by scripts/make-chapters.py.  Do not edit: change the order of the
% index and rerun the script (./preview-site.sh and CI do this).
%
% Usage in a sheet, after \newcounter{chap}:
%   \input{chapters}\chapnumber{<sheet ID>}
\makeatletter
\providecommand\chapters@set[2]{\@namedef{chapters@#1}{#2}}
\chapters@set{20261125T001000}{1}% zeros-to-ramp-sums
\chapters@set{20261125T001100}{2}% inert-primes-and-estimate-b
\chapters@set{20261125T005000}{3}% galois-fourier-and-averaging
\chapters@set{20261125T005500}{4}% columns-from-a-sphere
\chapters@set{20261125T000000}{5}% slope-method
\chapters@set{20261126T030100}{6}% a-twist-for-each-prime
\chapters@set{20261125T003000}{7}% the-chain-and-the-square-root
\chapters@set{20261125T010000}{8}% rank-lemma-and-zero-estimate
\chapters@set{20261125T010500}{9}% poincare-series-and-near-orthogonality
\chapters@set{20261124T223000}{10}% elliptic-curves-for-the-proof
\chapters@set{20261125T009500}{11}% fock-model-and-theta-intertwiners
\chapters@set{20261124T223100}{12}% complex-multiplication-and-inert-reduction
\chapters@set{20261125T008000}{13}% one-cm-determinant
\chapters@set{20261124T214500}{14}% supersingular-theta-divisibility
\chapters@set{20261124T214600}{15}% the-isotropic-plane
\chapters@set{20261125T013700}{16}% pairings-on-a-formal-group-of-height-two
\chapters@set{20261125T013500}{17}% hecke-nondegeneracy-on-the-tree
\chapters@set{20261126T030200}{18}% the-bound-is-attained
\chapters@set{20261125T009000}{19}% kronecker-limit-and-chowla-selberg
\chapters@set{20261125T002000}{20}% definite-theta-lifts
\chapters@set{20261125T015500}{21}% siegel-weil-in-the-smallest-case
\chapters@set{20261125T002100}{22}% definite-theta-lifts-and-norms
\chapters@set{20261125T016500}{23}% hecke-operators-and-theta-lifts
\chapters@set{20261125T016700}{24}% divisibility-along-the-chain
\chapters@set{20261125T004000}{25}% denominators-and-sha
\chapters@set{20261125T004100}{26}% tate-twists-and-sha
\chapters@set{20261126T020430}{27}% the-minus-part-by-gauss-sums
\chapters@set{20261125T014100}{28}% congruence-modules
\chapters@set{20261125T014500}{29}% eisenstein-congruences-and-selmer-classes
\chapters@set{20261125T014700}{30}% pseudocharacters-and-the-length-bound
\chapters@set{20261125T011000}{31}% galois-representations-black-box
\chapters@set{20261126T030300}{32}% the-data
\chapters@set{20261126T020210}{33}% bessel-poincare-series-and-their-gram-operators
\chapters@set{20261126T030400}{34}% the-seesaw-identity
\chapters@set{20261126T030500}{35}% divisibility-of-the-composite-lift
\chapters@set{20261126T030600}{36}% theorem-a-short-forms
\chapters@set{20261126T030700}{37}% the-determinant-of-a-chain
\chapters@set{20261126T030800}{38}% the-cost-of-one-lift
\chapters@set{20261126T030900}{39}% denominators-of-projections-to-lifts
\chapters@set{20261126T031000}{40}% twists-with-few-congruences
\chapters@set{20261126T031100}{41}% theorem-c-and-the-main-inequality
\providecommand\chapnumber[1]{%
  \@ifundefined{chapters@#1}%
    {\PackageError{chapters}{Sheet #1 is not listed in the index}%
      {List it in 20261125T012000--grh-workbook__guide.org\MessageBreak
       and run scripts/make-chapters.py.}}%
    {\edef\thechap{\@nameuse{chapters@#1}}}}
\makeatother
\chapnumber{<sheet ID>}
\makeatletter
\providecommand\chapters@set[2]{\@namedef{chapters@#1}{#2}}
\chapters@set{20261125T001000}{1}% zeros-to-ramp-sums
\chapters@set{20261125T001100}{2}% inert-primes-and-estimate-b
\chapters@set{20261125T005000}{3}% galois-fourier-and-averaging
\chapters@set{20261125T005500}{4}% columns-from-a-sphere
\chapters@set{20261125T000000}{5}% slope-method
\chapters@set{20261126T030100}{6}% a-twist-for-each-prime
\chapters@set{20261125T003000}{7}% the-chain-and-the-square-root
\chapters@set{20261125T010000}{8}% rank-lemma-and-zero-estimate
\chapters@set{20261125T010500}{9}% poincare-series-and-near-orthogonality
\chapters@set{20261124T223000}{10}% elliptic-curves-for-the-proof
\chapters@set{20261125T009500}{11}% fock-model-and-theta-intertwiners
\chapters@set{20261124T223100}{12}% complex-multiplication-and-inert-reduction
\chapters@set{20261125T008000}{13}% one-cm-determinant
\chapters@set{20261124T214500}{14}% supersingular-theta-divisibility
\chapters@set{20261124T214600}{15}% the-isotropic-plane
\chapters@set{20261125T013700}{16}% pairings-on-a-formal-group-of-height-two
\chapters@set{20261125T013500}{17}% hecke-nondegeneracy-on-the-tree
\chapters@set{20261126T030200}{18}% the-bound-is-attained
\chapters@set{20261125T009000}{19}% kronecker-limit-and-chowla-selberg
\chapters@set{20261125T002000}{20}% definite-theta-lifts
\chapters@set{20261125T015500}{21}% siegel-weil-in-the-smallest-case
\chapters@set{20261125T002100}{22}% definite-theta-lifts-and-norms
\chapters@set{20261125T016500}{23}% hecke-operators-and-theta-lifts
\chapters@set{20261125T016700}{24}% divisibility-along-the-chain
\chapters@set{20261125T004000}{25}% denominators-and-sha
\chapters@set{20261125T004100}{26}% tate-twists-and-sha
\chapters@set{20261126T020430}{27}% the-minus-part-by-gauss-sums
\chapters@set{20261125T014100}{28}% congruence-modules
\chapters@set{20261125T014500}{29}% eisenstein-congruences-and-selmer-classes
\chapters@set{20261125T014700}{30}% pseudocharacters-and-the-length-bound
\chapters@set{20261125T011000}{31}% galois-representations-black-box
\chapters@set{20261126T030300}{32}% the-data
\chapters@set{20261126T020210}{33}% bessel-poincare-series-and-their-gram-operators
\chapters@set{20261126T030400}{34}% the-seesaw-identity
\chapters@set{20261126T030500}{35}% divisibility-of-the-composite-lift
\chapters@set{20261126T030600}{36}% theorem-a-short-forms
\chapters@set{20261126T030700}{37}% the-determinant-of-a-chain
\chapters@set{20261126T030800}{38}% the-cost-of-one-lift
\chapters@set{20261126T030900}{39}% denominators-of-projections-to-lifts
\chapters@set{20261126T031000}{40}% twists-with-few-congruences
\chapters@set{20261126T031100}{41}% theorem-c-and-the-main-inequality
\providecommand\chapnumber[1]{%
  \@ifundefined{chapters@#1}%
    {\PackageError{chapters}{Sheet #1 is not listed in the index}%
      {List it in 20261125T012000--grh-workbook__guide.org\MessageBreak
       and run scripts/make-chapters.py.}}%
    {\edef\thechap{\@nameuse{chapters@#1}}}}
\makeatother
\chapnumber{<sheet ID>}
\makeatletter
\providecommand\chapters@set[2]{\@namedef{chapters@#1}{#2}}
\chapters@set{20261125T001000}{1}% zeros-to-ramp-sums
\chapters@set{20261125T001100}{2}% inert-primes-and-estimate-b
\chapters@set{20261125T005000}{3}% galois-fourier-and-averaging
\chapters@set{20261125T005500}{4}% columns-from-a-sphere
\chapters@set{20261125T000000}{5}% slope-method
\chapters@set{20261126T030100}{6}% a-twist-for-each-prime
\chapters@set{20261125T003000}{7}% the-chain-and-the-square-root
\chapters@set{20261125T010000}{8}% rank-lemma-and-zero-estimate
\chapters@set{20261125T010500}{9}% poincare-series-and-near-orthogonality
\chapters@set{20261124T223000}{10}% elliptic-curves-for-the-proof
\chapters@set{20261125T009500}{11}% fock-model-and-theta-intertwiners
\chapters@set{20261124T223100}{12}% complex-multiplication-and-inert-reduction
\chapters@set{20261125T008000}{13}% one-cm-determinant
\chapters@set{20261124T214500}{14}% supersingular-theta-divisibility
\chapters@set{20261124T214600}{15}% the-isotropic-plane
\chapters@set{20261125T013700}{16}% pairings-on-a-formal-group-of-height-two
\chapters@set{20261125T013500}{17}% hecke-nondegeneracy-on-the-tree
\chapters@set{20261126T030200}{18}% the-bound-is-attained
\chapters@set{20261125T009000}{19}% kronecker-limit-and-chowla-selberg
\chapters@set{20261125T002000}{20}% definite-theta-lifts
\chapters@set{20261125T015500}{21}% siegel-weil-in-the-smallest-case
\chapters@set{20261125T002100}{22}% definite-theta-lifts-and-norms
\chapters@set{20261125T016500}{23}% hecke-operators-and-theta-lifts
\chapters@set{20261125T016700}{24}% divisibility-along-the-chain
\chapters@set{20261125T004000}{25}% denominators-and-sha
\chapters@set{20261125T004100}{26}% tate-twists-and-sha
\chapters@set{20261126T020430}{27}% the-minus-part-by-gauss-sums
\chapters@set{20261125T014100}{28}% congruence-modules
\chapters@set{20261125T014500}{29}% eisenstein-congruences-and-selmer-classes
\chapters@set{20261125T014700}{30}% pseudocharacters-and-the-length-bound
\chapters@set{20261125T011000}{31}% galois-representations-black-box
\chapters@set{20261126T030300}{32}% the-data
\chapters@set{20261126T020210}{33}% bessel-poincare-series-and-their-gram-operators
\chapters@set{20261126T030400}{34}% the-seesaw-identity
\chapters@set{20261126T030500}{35}% divisibility-of-the-composite-lift
\chapters@set{20261126T030600}{36}% theorem-a-short-forms
\chapters@set{20261126T030700}{37}% the-determinant-of-a-chain
\chapters@set{20261126T030800}{38}% the-cost-of-one-lift
\chapters@set{20261126T030900}{39}% denominators-of-projections-to-lifts
\chapters@set{20261126T031000}{40}% twists-with-few-congruences
\chapters@set{20261126T031100}{41}% theorem-c-and-the-main-inequality
\providecommand\chapnumber[1]{%
  \@ifundefined{chapters@#1}%
    {\PackageError{chapters}{Sheet #1 is not listed in the index}%
      {List it in 20261125T012000--grh-workbook__guide.org\MessageBreak
       and run scripts/make-chapters.py.}}%
    {\edef\thechap{\@nameuse{chapters@#1}}}}
\makeatother
\chapnumber{20261126T080500}
\renewcommand{\thesection}{\thechap.\arabic{section}}

\theoremstyle{plain}
\newtheorem{theorem}{Theorem}[section]
\theoremstyle{definition}
\newtheorem{definition}[theorem]{Definition}
\newtheorem{problem}{Problem}[section]
\theoremstyle{remark}
\newtheorem*{leanremark}{In the Lean formalization}

\newcommand{\F}{\mathbb{F}}
\newcommand{\R}{\mathbb{R}}
\newcommand{\Z}{\mathbb{Z}}
\newcommand{\GL}{\mathrm{GL}}
\newcommand{\SL}{\mathrm{SL}}
\newcommand{\xv}{\mathsf{x}}
\newcommand{\yv}{\mathsf{y}}
\newcommand{\Tv}{\mathsf{T}}
\newcommand{\Sv}{\mathsf{S}}
\newcommand{\wt}{\mathrm{wt}}
\newcommand{\zo}{\mathrm{zo}}
\newcommand{\po}{\mathrm{po}}
\newcommand{\aux}{\mathrm{aux}}
\newcommand{\Deep}{\mathrm{Deep}}
\newcommand{\tl}[1]{\texttt{#1}}
\setlength{\emergencystretch}{2em}

\title{Least terms and Pieri's inequalities}
\author{}
\date{November 26, 2026}

\begin{document}
\maketitle
\refstepcounter{chap}\label{chap}

For small exceptional primes, the proof of Theorem~A(b) uses a near-surjectivity bound for scaled Hecke operators modulo prime powers.  Chapter~\ref{d4-chap}, \emph{A Hecke operator is nearly onto} (D4 below), realizes these operators on edges.  For suitable variable weights, the root lemma bounds the defect of a term attaining a finite least order across translates by determinant-one matrices acting in a fixed ordered pair of rows.  Order combines monomial weight with coefficient divisibility.  The proof counts zeros of an auxiliary polynomial with multiplicity in the style of Stepanov \cite[\S I]{Bo}.  Pieri's inequalities say that the defect, a linear expression in the exponents, is nonpositive for each set of variables, called a \emph{letter}.  For one letter, the fraction of the exponent simplex allowed is the reciprocal of the number of full flags, chains of subspaces with one of each dimension.  D4 uses the corresponding volume fraction to bound the dimension of the space of forms with bounded defects.  A related zero count bounds the sum of the dimensions of the spaces of deep forms, whose translates have only terms of sufficiently large total exponent in designated variables (Section~\ref{d4-sec:ends}).  The exercises establish Pieri's inequalities for least terms in the one-letter setting and in D4's two-letter model over \(\F_{25}\).

\paragraph{Prerequisites.}
Multiplicity of a zero of a binary form on the projective line.  Chapter~\ref{d4-chap} gives the matrix action on several sets of variables by means of a Frobenius cycle (Definition~\ref{d4-def:cycle}), the formula for the defect (Definition~\ref{d4-def:defect}), and the model used here over \(\F_{25}\).

We start with binary forms to see how a root count can restrict the least exponent under linear changes of variables.

\section{One letter}\label{sec:one-letter}

Let \(F\ne0\) be a homogeneous polynomial of degree \(\mu\) in two variables \(\xv,\yv\) over \(\F_q\).  For a matrix \(g=(g_{ij})\), replace variable \(j\) by the linear combination of the variables with coefficients \(g_{ij}\).  This substitution is the \emph{column action}, and the resulting polynomial \(g\cdot F\) is a translate of \(F\).  The first problems bound the least exponent of \(\yv\) across translates and give a form attaining the bound.

\begin{problem}\label{pr:multiplicities}
The sum of the multiplicities of \(F\) at the \(q+1\) rational points of \(\mathbb{P}^1\) is at most \(\mu\).  At some rational point the multiplicity is at most \(\mu/(q+1)\).
\end{problem}

\begin{problem}\label{pr:one-letter}
As \(g\) runs through \(\SL_2(\F_q)\), the least exponent of \(\yv\) in \(g\cdot F\) runs through the multiplicities at the rational points.  If \(\xv^{\mu-b}\yv^b\) occurs in some translate and \(b\) is least among all terms of all translates, then \(\mu-b\ge qb\).
\end{problem}

\begin{problem}\label{pr:invariant-one}
The form
\[
  F=\xv^q\yv-\xv\yv^q=\yv\prod_{v\in\F_q}(\xv-v\yv)
\]
is fixed by \(\SL_2(\F_q)\) and attains equality in Problem~\ref{pr:one-letter}.  For \(q=3\) its square is \(\xv^6\yv^2+\xv^4\yv^4+\xv^2\yv^6\).
\end{problem}

\begin{problem}\label{pr:dehomogenize}
Take a translate \(g\cdot F\) in which the minimum \(b\) of Problem~\ref{pr:one-letter} occurs.  The polynomial \((g\cdot F)(\Tv,1)\) has degree \(\mu-b\) and is divisible by \((\Tv-v)^b\) for every \(v\in\F_q\).  In particular \(qb\le\mu-b\).
\end{problem}

We now bound the least exponents in more variables and use the bound to count coefficients that determine a form.  Let \(n\ge1\) be an integer and let \(F\) be a nonzero form of degree \(\mu\) in \(z_0,\dots,z_{n-1}\) over \(\F_q\), with the same column action.  Row \(j\) means the index of \(z_j\).  Compare exponent vectors \(\alpha\) by reading \(\alpha_{n-1},\alpha_{n-2},\dots,\alpha_0\) in order, with the smaller entry at the first difference coming first.  Write \(z^\alpha\) for the monomial with exponents \(\alpha\) and \([z^\alpha]F\) for its coefficient.

\begin{problem}\label{pr:pieri-one}
A least exponent among all terms of all \(\GL_n(\F_q)\)-translates satisfies the one-letter Pieri inequalities \(\alpha_j\ge q\alpha_{j+1}\) for \(0\le j<n-1\).
\end{problem}

\begin{problem}\label{pr:triangular}
Under a lower triangular invertible matrix, the least exponent of a nonzero form stays the same, and its coefficient is multiplied by a nonzero scalar.
\end{problem}

\begin{problem}\label{pr:flags}
The number of full flags in \(\F_q^n\) is \(\prod_{j=1}^n\frac{q^j-1}{q-1}\).  The same number counts the right cosets of the lower triangular subgroup in \(\GL_n(\F_q)\).
\end{problem}

\begin{problem}\label{pr:flag-count}
The coefficients at exponents satisfying Problem~\ref{pr:pieri-one}, in translates by representatives of the right cosets in Problem~\ref{pr:flags}, determine \(F\).  Hence \(\binom{\mu+n-1}{n-1}\) is at most the number of flags times the number of exponents of degree \(\mu\) satisfying Problem~\ref{pr:pieri-one}.  For \((q,n,\mu)=(3,2,4),(2,3,7),(3,3,5)\), the triples (dimension, exponent count, flag count) are \((5,2,4),(36,4,21),(21,2,52)\), respectively.
\end{problem}

\begin{problem}\label{pr:cone-share}
The part of the simplex \(\alpha\ge0,\ \sum_j\alpha_j\le1\) satisfying the inequalities of Problem~\ref{pr:pieri-one} has the volume share
\[
  \prod_{j=1}^n\frac{q-1}{q^j-1}.
\]
The volume share is \(1/(q+1)\) for \(n=2\) and \(1/52\) for \(n=3,\ q=3\).  For \(n=2\) the count in Problem~\ref{pr:flag-count} is at most the dimension plus \(q\):
\[
  (q+1)\bigl(\lfloor \mu/(q+1)\rfloor+1\bigr)-(\mu+1)\le q.
\]
\end{problem}

With two sets of variables, the next bound will involve the exponents in both sets at once.

\section{Two letters over \texorpdfstring{\(\F_{25}\)}{F25}}\label{sec:two-letters}

In the action of Definition~\ref{d4-def:cycle}, the map \(\iota_s\) gives the matrix coefficients for letter \(s\), and the row index specifies a variable within the letter.  We continue to write \(z^\alpha\) and \([z^\alpha]F\) for a monomial and its coefficient.

\begin{definition}\label{def:root-variables}
Fix an ordered pair of distinct rows \(a,a'\).  In an even letter \(s\), \(\xv_s,\yv_s\) have rows \(a,a'\), respectively.  In an odd letter, the rows are \(a',a\), respectively.  To compute the defect \(d_i(\alpha)\) for this pair, use the rows of \(\yv_s,\xv_s\) as \(u_s,\bar u_s\) in Definition~\ref{d4-def:defect}.  A block is a group of consecutive rows.  \emph{Pieri's inequalities} on a block require \(d_i(\alpha)\le0\) at every letter \(i\) for each adjacent pair \(a,a+1\) in the block.
\end{definition}

We now enumerate the small-degree pairs that satisfy both inequalities in D4's model.  Reduce D4's model to \(\F_{25}\).  There are two letters and two rows (\(N_0=n=2\)), with \(P=5\), \(q=25\), and maps \(\iota_0(v)=v,\ \iota_1(v)=v^5\).  For rows \(0,1\), the variables \((\xv_0,\yv_0,\xv_1,\yv_1)\) are \((z_0,z_1,z'_1,z'_0)\) in D4.  A term of bidegree \((\mu,\mu)\) has total degree \(\mu\) in each letter.  Write \(b,c\) for its exponents on \(\yv_0,\yv_1\).  For the pair \((b,c)\), the defects of Problem~\ref{d4-pr:cone} are
\[
  24 d_0=26b+10c-6\mu,\qquad 24 d_1=10b+26c-6\mu.
\]

\begin{problem}\label{pr:cone}
The pairs \((b,c)\) with both defects nonpositive are \((0,0)\) alone for \(0\le\mu\le4\), the three pairs \((0,0),(0,1),(1,0)\) for \(\mu=5\), and those three together with \((1,1)\) for \(\mu=6\).
\end{problem}

To ask which pairs can occur as least terms, we need a way to compare terms by weight and to track those weights under substitution.

\begin{definition}\label{def:gauss-order}
Fix an element \(\varpi\) of a commutative coefficient ring.  The order \(\operatorname{ord}_{\varpi}(h)\) of a coefficient \(h\) is the supremum of the nonnegative integers \(j\) for which \(\varpi^j\) divides \(h\).  Assign a real weight \(\wt(z)\) to each variable \(z\).  The monomial \(z^\alpha\) has weight \(\wt\cdot\alpha:=\sum_z\wt(z)\alpha(z)\).  The \emph{Gauss order} \(\operatorname{ord}_{\wt} F\) is the least of \(\operatorname{ord}_{\varpi}([z^\alpha]F)+\wt\cdot\alpha\) over the terms of \(F\), and is \(+\infty\) for \(F=0\).  A bound \(\operatorname{ord}_{\wt} F\ge t\) requires each coefficient \([z^\alpha]F\) to be divisible by \(\varpi^{\max(0,\lceil t-\wt\cdot\alpha\rceil)}\).
\end{definition}

In the field model \(\varpi=0\), so the Gauss order is the least weight of a term with nonzero coefficient.  Use weights \(0\) on \(\xv_0,\xv_1\) and weights in \([1/2,1]\) on \(\yv_0,\yv_1\).  Adding the same weight to both variables of a letter adds a constant to the weight of every term of a fixed bidegree.  We find the least term of an invariant form and then give weight bounds for substitutions.

\begin{problem}\label{pr:invariant}
The invariant from Problem~\ref{d4-pr:invariant} is
\[
  (\xv_0^5\yv_1+\yv_0^5\xv_1)(\xv_0\yv_1^5+\yv_0\xv_1^5).
\]
Its pairs \((b,c)\) are \((0,6),(1,1),(5,5),(6,0)\).  The least term is \((1,1)\), where \(d_0=d_1=0\).
\end{problem}

\begin{problem}\label{pr:substitution}
Over the field, give real weights to the variables of a source and a target polynomial ring.  Suppose a substitution sends each source variable to a polynomial of Gauss order at least the source variable's weight.  Then every polynomial has Gauss order at least as large after substitution as before.  Suppose also that each variable's image is the sum of a main image and extra summands, with the main image meeting the same lower bound and each extra summand exceeding the source weight by at least \(\varepsilon>0\).  The difference between the polynomial obtained from all summands and the polynomial obtained from the main images has Gauss order at least the original order plus \(\varepsilon\).
\end{problem}

To bound the defects of a least term, we extract a polynomial in one variable whose zeros can be counted.

\section{The auxiliary polynomial}\label{sec:aux}

Let \(F\) be a form of bidegree \((\mu,\mu)\) in the field model.  Take \(z^\alpha\) with nonzero coefficient in \(F\) and assume that its weight is least among all terms of all \(\SL_2(\F_{25})\)-translates of \(F\).  We first treat the defect at letter \(0\).

\begin{definition}\label{def:aux}
The substitution \(\Sv\) in the two letters is given by
\[
\begin{aligned}
  \xv_0&\longmapsto\Tv^{25}\xv_0, & \yv_0&\longmapsto\xv_0,\\
  \xv_1&\longmapsto\Tv^{125}\xv_1-\Tv^5\yv_1,&
  \yv_1&\longmapsto\yv_1-\xv_1 .
\end{aligned}
\]
In \(\Sv(F)\), regard \(\Tv\) as a coefficient variable.  Let \(\aux_\alpha(\Tv)\) be the coefficient of the monomial obtained from \(z^\alpha\) by replacing \(\yv_0\) by \(\xv_0\).  Assign zero weights \(25,5\) to \(\yv_0,\yv_1\) and pole weights \(25,125\) to \(\xv_0,\xv_1\).  The unlisted entries of each list are zero.  Summing each list of weights with the exponents of \(z^\alpha\) gives
\[
  \zo(\alpha):=25b+5c,\qquad
  \po(\alpha):=25(\mu-b)+125(\mu-c).
\]
\end{definition}

The next problems relate the degree and zero multiplicities of \(\aux_\alpha\) to the defect.

\begin{problem}\label{pr:aux-degree}
Give \(\Tv\) weight \(-\wt(\yv_0)/25\), and raise the weight of each output variable by its pole weight times \(\wt(\yv_0)/25\).  In each substitution in Definition~\ref{def:aux}, the first summand then has the weight of the input variable, and any second summand exceeds the input weight by at least \(1/20\).  The degree of \(\aux_\alpha\) is \(\po(\alpha)\) and its leading coefficient is \([z^\alpha]F\).
\end{problem}

\begin{problem}\label{pr:aux-action}
For \(v\in\F_{25}\) let \(g_v=\left(\begin{smallmatrix}0&1\\-1&-v\end{smallmatrix}\right)\).  Under \(g_v\) followed by \(\Sv\) and then \(\Tv\mapsto\Tv+v\), the images of \(\xv_0,\yv_0,\xv_1,\yv_1\) are, in order,
\[
  -\xv_0,\qquad \Tv^{25}\xv_0,\qquad
  \yv_1-\xv_1,\qquad -\Tv^{125}\xv_1+\Tv^5\yv_1 .
\]
\end{problem}

\begin{problem}\label{pr:aux-zeros}
Give \(\Tv\) weight \(\wt(\yv_0)/25\), and lower the weight of each output variable by its zero weight times \(\wt(\yv_0)/25\).  The images in Problem~\ref{pr:aux-action} have Gauss order at least the original weights of the respective inputs.  Consequently \((\Tv-v)^{\zo(\alpha)}\) divides \(\aux_\alpha(\Tv)\) for every \(v\in\F_{25}\).
\end{problem}

\begin{problem}\label{pr:root-model}
The orders satisfy \(25\zo(\alpha)\le\po(\alpha)\), equivalently \(d_0(\alpha)\le0\).  Also \(d_1(\alpha)\le0\).
\end{problem}

\begin{problem}\label{pr:aux-invariant}
For the invariant of Problem~\ref{pr:invariant},
\[
  \Sv(F)=\mathsf{x}_0^6\,\mathsf{x}_1^5\,\mathsf{y}_1\,(\Tv^{25}-\Tv)^{30}.
\]
For \((b,c)=(1,1)\), the auxiliary polynomial has degree \(750=\po(\alpha)\) and multiplicity \(30=\zo(\alpha)\) at each point of \(\F_{25}\).
\end{problem}

In the field model, depth is \(b+c\), so \(\Deep_1\) is the space of forms whose \((0,0)\) coefficient vanishes in every \(\GL_2(\F_{25})\)-translate (Definition~\ref{d4-def:deep}).  The next problem computes the dimension of \(\Deep_1\) in small degrees and identifies the possible least terms in degree \(5\).

\begin{problem}\label{pr:deep-ranks}
(i) Verify by computer that for \(\mu=5,6\) the linear map recording the coefficients at the pairs in Problem~\ref{pr:cone} in each of the \(26\) translates from Problem~\ref{d4-pr:cycle-model} has rank \(36,49\), respectively.  The maps recording just the \((0,0)\) coefficient have rank \(26\) in both degrees.  Thus \(\dim\Deep_1=10,23\), as in Problem~\ref{d4-pr:deep-table}.

(ii) For \(0\le\mu\le4\), \(\Deep_1=0\).  If \(0\ne F\in\Deep_1\) and \(\mu=5\), a term attaining the least weight among all \(\SL_2(\F_{25})\)-translates has pair \((1,0)\) or \((0,1)\), with \((\po,\zo)\) equal to \((725,25)\) or \((625,5)\), respectively.
\end{problem}

The general root lemma extends the bound on defects to cycles with more variables and to coefficients in a ring.

\section{The root lemma}\label{sec:root}

\begin{theorem}\label{thm:root}
Let a Frobenius cycle of length \(N_0\ge1\) from \(W'\) to \(O\) be given as in Definition~\ref{d4-def:cycle}, allowing any commutative coefficient ring \(O\).  Fix distinct rows \(a,a'\) and real weights such that \(\wt(\yv_s)-\wt(\xv_s)\in[1/2,1]\) for every \(s\).  Let \(F\) be a polynomial over \(O\), \(\alpha\) an exponent, and \(l\ge0\) an integer.  Suppose \(\varpi^{l+1}\) does not divide \([z^\alpha]F\) and every translate by \(\SL_2(W')\) in rows and columns \(a,a'\) has Gauss order at least \(l+\wt\cdot\alpha\).  Then at every letter \(i\), with \([s-i]\) the least nonnegative residue modulo \(N_0\),
\[
  q\alpha(\yv_i)+\sum_{s\ne i}P^{[s-i]}\alpha(\yv_s)
  \ \le\ \sum_s P^{[s-i]}\alpha(\xv_s).
\]
Equivalently, \(d_i(\alpha)\le0\) for the ordered pair \(a,a'\).
\end{theorem}

Problems~\ref{pr:aux-degree}--\ref{pr:root-model} prove the theorem on forms of bidegree \((\mu,\mu)\) in the model with \(l=0\).

The general root lemma is taken as given.

Section~\ref{d4-sec:middle} of D4 divides the rows into two blocks of sizes \(r_0\) and \(n-r_0\).  The dimension estimate concerns forms whose translates by matrices acting separately on the two blocks have zero coefficient wherever the defect for the pair \((r_0-1,r_0)\) exceeds \(3t_0\) at some letter.  To bound the dimension, the proof packs a family of these forms into a single polynomial and selects a term in a block diagonal translate of that polynomial.  The term has nonpositive defect at every letter along every adjacent pair within either block.  The proof divides its exponent vector by the packing scale and then adds a shift.  It counts lattice points in the resulting degree range, with coordinates for all rows and letters and nonpositive defect at every letter along every adjacent pair, including the boundary pair.  The divisibility condition on the average in the same section uses the pair \((r_0-1,r_0)\) at the boundary between the blocks.

\begin{leanremark}
\raggedright
In \cite{GL}, Theorem~\ref{thm:root} is \tl{FrobeniusCycle.\allowbreak defect\_\allowbreak nonpos\_\allowbreak of\_\allowbreak extreme\_\allowbreak term} in \tl{RootLemma.lean}.  Definition~\ref{def:aux} specializes \tl{FrobeniusCycle.\allowbreak RootAt.\allowbreak mainSubst} minus \tl{crossSubst}, with coefficient \tl{mergedCoeff} in the same namespace.

The one-letter volume formula in Problem~\ref{pr:cone-share} is \tl{volume\_\allowbreak real\_\allowbreak pieriCone\_\allowbreak inter\_\allowbreak degreeSlab\_\allowbreak one} in \tl{PieriCone.lean}.  The volume fraction used in D4 is given by \tl{FrobeniusCycle.\allowbreak prod\_\allowbreak mul\_\allowbreak volume\_\allowbreak defectCone\_\allowbreak inter\_\allowbreak degreeSlab} in \tl{DefectCount.lean}.
\end{leanremark}

The hints give starting points for the binary-form and flag counts, then for the auxiliary polynomial and the dimension calculation over \(\F_{25}\).

\section*{Hints}

\subsection*{\S\ref{sec:one-letter}}
\begin{description}
\item[\ref{pr:one-letter}] Use the local coordinate at \([1:0]\).
\item[\ref{pr:invariant-one}] Use Frobenius on the two substituted linear forms.
\item[\ref{pr:dehomogenize}] Use Problem~\ref{pr:one-letter} in the chart \(\yv=1\).
\item[\ref{pr:pieri-one}] Use the subgroup \(\SL_2(\F_q)\) acting in rows \(j,j+1\).
\item[\ref{pr:flags}] Count flags by induction on the dimension.
\item[\ref{pr:flag-count}] Use Problem~\ref{pr:triangular}.
\item[\ref{pr:cone-share}] Use a triangular change of coordinates adapted to the inequalities.
\end{description}

\subsection*{\S\ref{sec:aux}}
\begin{description}
\item[\ref{pr:aux-degree}] Use Problem~\ref{pr:substitution} with the first summands as the main substitution.
\item[\ref{pr:aux-action}] Read columns of \(g_v\) and its fifth-power inverse transpose.  Use \(v^{25}=v\).
\item[\ref{pr:aux-zeros}] Apply Problem~\ref{pr:substitution} to a translate of \(F\).  Compare the output weights of \(\xv_0\) and \(\yv_0\).
\item[\ref{pr:root-model}] For \(d_0\), use Problem~\ref{d4-pr:one-variable}(i).  For \(d_1\), try renumbering the letters and rows.
\item[\ref{pr:aux-invariant}] Substitute in the two factors separately.
\item[\ref{pr:deep-ranks}] In SageMath, use \tl{monomial\_coefficient} and \tl{rank}.  For (i), write each matrix as a lower triangular matrix times one of the representatives.
\end{description}

\begin{thebibliography}{GL}
\bibitem[Bo]{Bo} E.~Bombieri, \emph{Counting points on curves over finite fields (d'apr\`es S.~A.~Stepanov)}, S\'eminaire Bourbaki, Exp.~430, Lecture Notes in Math. \textbf{383}, Springer, 1974, 234--241.
\bibitem[GL]{GL} L.~Sernau, \emph{GRH for Dirichlet \(L\)-functions, in Lean~4}, \href{https://github.com/lsernau/grh-lean/tree/72fead102f663d5c7a407964a13dc0cf9d977207/grh}{readable edition}, commit \texttt{72fead1}.
\end{thebibliography}

\end{document}
